\Needspace{10\baselineskip}
\begin{center}\emph{Die Reste $+2$ und $-2$.}\end{center}

\Article{112}
Wir gehen zu den Resten $+2$ und $-2$ über.

Sammeln wir aus Tabelle II alle Primzahlen, für die $+2$ ein Rest ist, so erhalten wir $7,17,23,31,41,47,71,73,79.89,97$. Man bemerkt leicht, dass keine davon die Form $8n+3$ oder $8n+5$ hat.

Sehen wir also nach, ob sich diese Beobachtung zur Gewissheit machen lässt.

Zunächst muss jede zusammengesetzte Zahl der Form $8n+3$ oder $8n+5$ einen Primfaktor einer der Formen $8n+3$ oder $8n+5$ enthalten. Denn aus Primzahlen der Formen $8n+1$, $8n+7$ allein lassen sich nur Zahlen der Form $8n+1$ oder $8n+7$ zusammensetzen. Gilt unsere Beobachtung also allgemein, so gibt es überhaupt keine Zahl der Form \SourceLeaf{90}$8n+3$, $8n+5$, für die $+2$ ein Rest ist. Jedenfalls gibt es unterhalb von $100$ keine solche Zahl, für die $+2$ ein Rest ist. Finden sich aber jenseits dieser Grenze solche Zahlen, so sei die kleinste $=t$. Dann hat $t$ entweder die Form $8n+3$ oder $8n+5$; $+2$ ist ein Rest davon, aber ein Nichtrest aller kleineren Zahlen dieser Formen. Wir setzen $2\equiv aa$ (mod.~$t$). Dabei lässt sich $a$ stets ungerade und $<t$ wählen (denn $a$ hat mindestens zwei positive Werte unterhalb von $t$, deren Summe $=t$ ist, sodass einer gerade und der andere ungerade ist; siehe Artikel 104 und 105). Dann sei $aa=2+tu$ oder $tu=aa-2$. Somit hat $aa$ die Form $8n+1$, also $tu$ die Form $8n-1$, und folglich $u$ die Form $8n+3$ oder $8n+5$, je nachdem, ob $t$ die zweite oder die erste Form hat. Aus der Gleichung $aa=2+tu$ folgt aber auch $2\equiv aa$ (mod.~$u$), \emph{das heißt,} $2$ ist auch ein Rest von $u$. Man sieht leicht, dass $u<t$ gilt. Also ist $t$ nicht die kleinste Zahl, die unserer Beobachtung widerspricht, entgegen der Voraussetzung. Damit gilt das durch Zahlenbeispiele gefundene Muster allgemein.

Verbinden wir dies mit der Aussage von Artikel 111, so erhalten wir folgende Sätze.

I. \emph{Für alle Primzahlen der Form $8n+3$ ist $+2$ ein Nichtrest, $-2$ dagegen ein Rest.}

II. \emph{Für alle Primzahlen der Form $8n+5$ sind sowohl $+2$ als auch $-2$ Nichtreste.}

\Article{113}
Ebenso finden wir durch Betrachtung der Tabelle II die folgenden Primzahlen, für die $-2$ ein Rest ist: $3,11,17,19,41,43,59,67,73,83,89,97$\HistoricalNote{*)}{Indem wir $-2$ als Produkt von $+2$ und $-1$ betrachten; siehe Artikel 111.}. Keine davon hat die Form $8n+5$, $8n+7$. Untersuchen wir also, ob sich auch diese Beobachtung als allgemeiner Satz beweisen lässt. Wie im vorigen Artikel zeigt man, dass jede zusammengesetzte Zahl der Form $8n+5$ oder $8n+7$ einen Primfaktor der Form $8n+5$ oder $8n+7$ enthält. Gilt unsere Beobachtung allgemein, so kann $-2$ überhaupt kein Rest einer Zahl der Form $8n+5$ oder $8n+7$ sein. Gäbe es solche Zahlen, so sei die kleinste $=t$, und wir setzen $-2=aa-tu$. Wählen wir wie oben $a$ ungerade und kleiner als $t$, so hat $u$ die Form $8n+5$ oder $8n+7$, je nachdem, ob $t$ die Form $8n+7$ oder $8n+5$ hat. Aus $aa+2=tu$ und $a<t$ folgt leicht, dass auch \SourceLeaf{91}$u$ kleiner als $t$ ist. Schließlich ist $-2$ auch ein Rest von $u$, \emph{das heißt,} $t$ ist nicht die kleinste Zahl, die unserer Beobachtung widerspricht, entgegen der Voraussetzung. Also muss $-2$ ein Nichtrest aller Zahlen der Formen $8n+5$, $8n+7$ sein.

\Needspace{6\baselineskip}
Mit der Aussage von Artikel 111 ergeben sich daraus folgende Sätze:

I.\quad \emph{Für alle Primzahlen der Form $8n+5$ sind sowohl $-2$ als auch $+2$ Nichtreste,} wie wir bereits im vorigen Artikel gefunden haben.

II.\quad \emph{Für alle Primzahlen der Form $8n+7$ ist $-2$ ein Nichtrest, $+2$ dagegen ein Rest.}

In beiden Beweisen hätten wir für $a$ auch einen geraden Wert wählen können. Dann hätten wir den Fall, dass $a$ die Form $4n+2$ hat, von dem Fall unterscheiden müssen, dass $a$ die Form $4n$ hat. Der weitere Beweis verläuft wie oben und bereitet keine Schwierigkeit.

\Article{114}
Ein Fall bleibt noch übrig: Die Primzahl hat die Form $8n+1$. Dieser Fall entzieht sich der vorigen Methode und erfordert ganz andere Kunstgriffe.

Für den Primzahlmodul $8n+1$ sei $a$ eine beliebige primitive Wurzel. Dann gilt (Artikel 62) $a^{4n}\equiv-1\,(\operatorname{mod.}8n+1)$. Diese Kongruenz lässt sich auch als $(a^{2n}+1)^2\equiv2a^{2n}\,(\operatorname{mod.}8n+1)$ oder als $(a^{2n}-1)^2\equiv-2a^{2n}$ schreiben. Daher sind sowohl $2a^{2n}$ als auch $-2a^{2n}$ Reste von $8n+1$. Da aber $a^{2n}$ ein Quadrat ist, das nicht durch den Modul teilbar ist, sind auch $+2$ und $-2$ Reste (Artikel 98).

\Article{115}
Es wird nützlich sein, noch einen anderen Beweis dieses Satzes anzufügen. Er steht zum vorigen in einer ähnlichen Beziehung wie der zweite Beweis (Artikel 109) des Satzes aus Artikel 108 zum ersten (Artikel 108). Mit diesen Fragen vertraute Leser werden dann leichter erkennen, dass weder das eine noch das andere Beweispaar so verschieden ist, wie es auf den ersten Blick scheinen mag.

\Needspace{5\baselineskip}
I.\quad Für einen beliebigen Primzahlmodul der Form $4m+1$ gibt es unter den kleineren Zahlen $1,2,3\ldots4m$ genau $m$, die einer vierten Potenz kongruent sein können; die übrigen $3m$ können es nicht.

Dies folgt leicht aus den Grundsätzen des vorigen Abschnitts, lässt sich aber auch ohne diese Grundsätze leicht beweisen. Wir haben nämlich gezeigt, dass für einen solchen Modul $-1$ \SourceLeaf{92}stets ein quadratischer Rest ist. Sei also $ff\equiv-1$. Ist $z$ eine beliebige nicht durch den Modul teilbare Zahl, so sind die vierten Potenzen der vier Zahlen $+z,-z,+fz,-fz$ (die offensichtlich inkongruent sind) untereinander kongruent. Die vierte Potenz einer Zahl, die keiner dieser vier kongruent ist, kann ihren vierten Potenzen nicht kongruent sein. Sonst hätte die Kongruenz $x^4\equiv z^4$, die vom vierten Grad ist, mehr als $4$ Wurzeln, entgegen Artikel 43. Somit liefern alle Zahlen $1,2,3,\ldots4m$ nur $m$ inkongruente vierte Potenzen. Unter denselben Zahlen sind $m$ diesen kongruent; die übrigen können keiner vierten Potenz kongruent sein.

\Needspace{5\baselineskip}
II.\quad Nach einem Primzahlmodul der Form $8n+1$ kann $-1$ einer vierten Potenz kongruent sein ($-1$ ist ein \emph{biquadratischer Rest} dieser Primzahl).

Denn die Anzahl aller biquadratischen Reste unterhalb von $8n+1$ (null ausgeschlossen) ist $=2n$, \emph{das heißt,} gerade. Ferner zeigt man leicht: Ist $r$ ein biquadratischer Rest von $8n+1$, so ist auch der Wert des Ausdrucks $\frac{1}{r}\,(\operatorname{mod.}8n+1)$ ein solcher Rest. Daher lassen sich alle biquadratischen Reste so in Klassen einteilen, wie wir in Artikel 109 die quadratischen Reste eingeteilt haben. Auch der übrige Teil des Beweises verläuft genau wie dort.

\Needspace{7\baselineskip}
III.\quad Nun sei $g^4\equiv-1$, und $h$ sei der Wert des Ausdrucks $\frac{1}{g}\,(\operatorname{mod.}8n+1)$. Dann gilt
\[
(g\pm h)^2=g^2+h^2\pm2gh\equiv g^2+h^2\pm2
\]
(wegen $gh\equiv1$). Aber es gilt $g^4\equiv-1$, also $-h^2\equiv g^4h^2\equiv g^2$, und schließlich $g^2+h^2\equiv0$ sowie $(g\pm h)^2\equiv\pm2$; \emph{das heißt,} sowohl $+2$ als auch $-2$ sind quadratische Reste von $8n+1$. \mbox{\emph{Q. E. D.}}

\Article{116}
Aus dem Vorigen ergibt sich außerdem leicht die allgemeine Regel: $+2$ \emph{ist ein Rest jeder Zahl, die weder durch $4$ noch durch eine Primzahl der Form $8n+3$ oder $8n+5$ teilbar ist; dagegen ist es ein Nichtrest aller übrigen Zahlen} (\emph{zum Beispiel} aller Zahlen der Formen $8n+3$, $8n+5$, ob sie Primzahlen oder zusammengesetzt sind).

$-2$ \emph{ist ein Rest jeder Zahl, die weder durch $4$ noch durch eine Primzahl der Form $8n+5$ oder $8n+7$ teilbar ist, und ein Nichtrest aller übrigen Zahlen.}

Diese eleganten Sätze waren schon dem scharfsinnigen Fermat bekannt, \emph{Op. Mathem., S.} 168. Einen Beweis, den er nach eigener Aussage besaß, hat er aber nirgends \SourceLeaf{93}mitgeteilt. Später suchte der bedeutende Euler stets vergeblich danach; der bedeutende Lagrange fand als Erster einen strengen Beweis, \emph{Nouv. M\'em. de l'Ac. de Berlin}, 1775, \emph{S.} 349, 351. Dies scheint Euler noch nicht bekannt gewesen zu sein, als er die in \emph{Opusc. Analyt.}, Bd. I, S. 259, enthaltene Abhandlung schrieb.

\Needspace{10\baselineskip}
\begin{center}\emph{Die Reste $+3$ und $-3$.}\end{center}

\Article{117}
Wir gehen zu den Resten $+3$ und $-3$ über und beginnen mit dem letzteren.

Aus Tabelle II finden wir folgende Primzahlen, für die $-3$ ein Rest ist: 3, 7. 13, 19, 31, 37, 43, 61, 67, 73, 79, 97. Keine davon hat die Form $6n+5$. Dass es auch jenseits der Tabelle keine Primzahlen dieser Form gibt, für die $-3$ ein Rest ist, beweisen wir so: Zunächst enthält jede zusammengesetzte Zahl der Form $6n+5$ einen Primfaktor derselben Form. Solange es also keine Primzahlen der Form $6n+5$ gibt, für die $-3$ ein Rest ist, gibt es auch keine solchen zusammengesetzten Zahlen. Gäbe es aber jenseits der Tabelle solche Zahlen, so sei die kleinste $=t$, und wir setzen $-3=aa-tu$. Wählen wir $a$ gerade und kleiner als $t$, so gilt $u<t$, und $-3$ ist ein Rest von $u$. Hat $a$ die Form $6n\pm2$, so hat $tu$ die Form $6n+1$ und damit $u$ die Form $6n+5$. \mbox{\emph{Q. E. A.}} Denn wir haben angenommen, dass $t$ die kleinste unserer Beobachtung widersprechende Zahl ist. Hat dagegen $a$ die Form $6n$, so hat $tu$ die Form $36n+3$, also $\frac{1}{3}tu$ die Form $12n+1$, und daher $\frac{1}{3}u$ die Form $6n+5$. Aber $-3$ ist auch ein Rest von $\frac{1}{3}u$, und es gilt $\frac{1}{3}u<t$, \mbox{\emph{Q. E. A.}} Somit kann $-3$ kein Rest einer Zahl der Form $6n+5$ sein.

Jede Zahl der Form $6n+5$ hat entweder die Form $12n+5$ oder $12n+11$. Die erste Form ist in $4n+1$ enthalten, die zweite in $4n+3$. Daher gelten die folgenden Sätze:

I.\quad \emph{Für jede Primzahl der Form $12n+5$ sind sowohl $-3$ als auch $+3$ Nichtreste.}

II.\quad \emph{Für jede Primzahl der Form $12n+11$ ist $-3$ ein Nichtrest, $+3$ dagegen ein Rest.}

\SourceLeaf{94}\Needspace{7\baselineskip}
\Article{118}
Aus Tabelle II finden wir die folgenden Zahlen, für die $+3$ ein Rest ist: 3, 11, 13, 23, 37, 47, 59, 61, 71, 73, 83, 97. Keine davon hat die Form $12n+5$ oder $12n+7$. Dass es überhaupt keine Zahlen der Formen $12n+5$, $12n+7$ gibt, für die $+3$ ein Rest ist, lässt sich genau wie in den Artikeln 112, 113 und 117 beweisen; deshalb verzichten wir hier darauf. Mit Artikel 111 erhalten wir die Sätze:

I.\quad \emph{Für jede Primzahl der Form $12n+5$ sind sowohl $+3$ als auch $-3$ Nichtreste} (wie schon im vorigen Artikel gefunden).

II.\quad \emph{Für jede Primzahl der Form $12n+7$ ist $+3$ ein Nichtrest, $-3$ dagegen ein Rest.}

\Article{119}
Für Zahlen der Form $12n+1$ liefert diese Methode aber kein Ergebnis; sie erfordern also besondere Kunstgriffe. Zahlenbeispiele legen nahe, dass $+3$ und $-3$ Reste aller Primzahlen dieser Form sind. Man muss nur beweisen, dass $-3$ ein Rest solcher Zahlen ist, denn dann muss auch $+3$ ein Rest sein (Artikel 111). Wir zeigen aber allgemeiner, dass $-3$ ein Rest jeder Primzahl der Form $3n+1$ ist.

Sei $p$ eine solche Primzahl und $a$ eine Zahl, die für den Modul $p$ zum Exponenten $3$ gehört (solche Zahlen gibt es nach Artikel 54, da $3$ ein Teiler von $p-1$ ist). Dann gilt $a^3\equiv1\,(\operatorname{mod.}p)$, \emph{das heißt,} $a^3-1$ oder $(a^2+a+1)(a-1)$ ist durch $p$ teilbar. Aber $a$ kann nicht $\equiv1\,(\operatorname{mod.}p)$ sein, da $1$ zum Exponenten $1$ gehört. Also ist $a-1$ nicht durch $p$ teilbar, wohl aber $a^2+a+1$ und daher auch $4aa+4a+4$; \emph{das heißt,} es gilt $(2a+1)^2\equiv-3\,(\operatorname{mod.}p)$, oder $-3$ ist ein Rest von $p$. \mbox{\emph{Q. E. D.}}

Dieser Beweis, der von den vorigen unabhängig ist, umfasst auch die Primzahlen der Form $12n+7$, die wir schon im vorigen Artikel behandelt haben.

Man kann noch bemerken, dass sich dieser Gedankengang auch nach dem Vorbild der Methode in den Artikeln 109 und 115 darstellen lässt. Der Kürze wegen gehen wir darauf nicht weiter ein.

\SourceLeaf{95}\Article{120}
Aus dem Vorigen ergeben sich leicht die folgenden Sätze (siehe Artikel 102, 103 und 105).

I. $-3$ \emph{ist ein Rest aller Zahlen, die weder durch $8$ noch durch $9$ noch durch eine Primzahl der Form $6n+5$ teilbar sind, und ein Nichtrest aller übrigen Zahlen.}

II. $+3$ \emph{ist ein Rest aller Zahlen, die weder durch $4$ noch durch $9$ noch durch eine Primzahl der Form $12n+5$ oder $12n+7$ teilbar sind, und ein Nichtrest aller übrigen Zahlen.}

Besonders sollte man sich diesen Spezialfall merken:

$-3$ ist ein \emph{Rest} aller Primzahlen der Form $3n+1$, oder gleichbedeutend \emph{aller Primzahlen, die Reste von $3$ sind}; dagegen ein \emph{Nichtrest} aller Primzahlen der Form $6n+5$ oder, wenn man $2$ ausschließt, aller Primzahlen der Form $3n+2$, \emph{das heißt aller, die Nichtreste von $3$ sind.} Alle übrigen Fälle folgen daraus offensichtlich von selbst.

\Needspace{10\baselineskip}
Die Aussagen über die Reste $+3$ und $-3$ waren schon Fermat bekannt, \emph{Opera Wallisii}, Bd. II, \emph{S.} 857. Doch der bedeutende Euler gab als Erster Beweise, \emph{Comm. nov. Petr.}, Bd. VIII, \emph{S.} 105 ff. Umso überraschender ist es, dass ihm die Beweise der Aussagen über die Reste $+2$ und $-2$, die auf ganz ähnlichen Kunstgriffen beruhen, trotz seines Scharfsinns stets entgingen. Siehe auch die Abhandlung des bedeutenden Lagrange, \emph{Nouv. M\'em. de l'Ac. de Berlin}, 1775, \emph{S.} 352.
